% =====================================================================
% The Cognitive Circuit Grimoire — front matter
% Noetic Compass Press · 2026
% =====================================================================

\frontmatter
\pagestyle{plainbook}

% --- Half-title --------------------------------------------------------
\thispagestyle{empty}
\begin{center}
  \vspace*{2.5in}
  {\huge\bfseries The Cognitive Circuit Grimoire\par}
\end{center}
% \cleardoublepage, not \clearpage: the title page opens on a recto and the
% copyright sits on the verso behind it.
\cleardoublepage

% --- Title page --------------------------------------------------------
\thispagestyle{empty}
\begin{center}
  \vspace*{1.1in}
  {\Huge\bfseries The Cognitive\\[0.10in] Circuit Grimoire\par}
  \vspace{0.55in}
  {\Large\itshape Spells, Hyperstitions, and Overton Engineering\\[0.10in]
   as Formal Cognitive Circuits\par}
  \vspace{0.35in}
  {\large Nine Boards, Drawn for Detection and for Use\par}
  \vspace{1.20in}
  {\Large Little Henderson\par}
  \vfill
  {\normalsize Noetic Compass Press\par}
  {\small 2026\par}
\end{center}
\clearpage

% --- Copyright ---------------------------------------------------------
\thispagestyle{empty}
\vspace*{\fill}
\begin{center}
\begin{minipage}{0.85\textwidth}
\small
\noindent
\textbf{The Cognitive Circuit Grimoire}\\
\textit{Spells, Hyperstitions, and Overton Engineering as Formal Cognitive
Circuits}\\[6pt]
Copyright \textcopyright{} 2026 Little Henderson. All rights reserved.\\[6pt]
No part of this publication may be reproduced, distributed, or transmitted in
any form or by any means, including photocopying, recording, or other electronic
or mechanical methods, without the prior written permission of the author, except
in the case of brief quotations embodied in critical reviews and certain other
noncommercial uses permitted by copyright law.\\[6pt]
First Edition, 2026. Published by Noetic Compass Press.\\[6pt]
Companion volumes: \textit{Thinking in Cognitive Architecture},
\textit{CAS PCBA Engineer}, and \textit{Reality Engineering: The Tootra
Kabbalistic System}.\\[6pt]
Framework: Cognitive Architecture System (CAS) v7.7.\\[6pt]
Base-framework intellectual property: USPTO Provisional Application
\#64/037,173. The Cognitive Junction Transistor (CJT), the Cognitive Circuit
Description Language (CCDL), and the Cognitive Circuit Board (CCB) reference
designs in this volume extend that foundation.\\[6pt]
ISBN: [pending publisher assignment]\\[6pt]
\textit{Product and company names mentioned in this book may be trademarks of
their respective owners, and are used here for identification and commentary
only.}
\end{minipage}
\end{center}
\vspace*{\fill}
\cleardoublepage

% --- Notice ------------------------------------------------------------
% This page is load-bearing and deliberately sits before the epigraph, where a
% reader opening the book in a shop will meet it. Several chapters analyse
% systems built to act on people who did not consent; the framing is structural
% analysis for detection and reversal, and it needs saying at the front in the
% plainest available language rather than only in the preface.
\thispagestyle{empty}
\vspace*{1.5in}
\begin{center}
\begin{minipage}{0.80\textwidth}
\small
\begin{center}\textbf{A NOTICE, BEFORE ANYTHING ELSE}\end{center}
\vspace{6pt}
\noindent
Five chapters of this book analyse systems built to act on people who did not
agree to it: conditioning, hypnosis applied by an operator, advertising, norm
progression, and information operations.

\vspace{6pt}
\noindent
They are written as structural analysis \emph{for detection, resistance, and
reversal}. The book explains how such a circuit is architected so that a reader
can recognise one; it does not supply procedures, scripts, or sequences for
running one on another person. Schematics whose purpose is to act on a
non-consenting person carry a \textsc{detection-only} banner, and every coercive
board in the catalogue ends with its counter-circuit.

\vspace{6pt}
\noindent
Work directed at the operator's own mind, with the operator's own consent, is a
different matter and is written fully and practically.

\vspace{6pt}
\noindent
\textbf{Nothing in this book is medical or psychiatric advice.} It diagnoses
nothing and treats nothing. Terms such as \emph{dissociation}, \emph{compulsion},
and \emph{despair} are used here to name circuit states, not clinical
conditions. If you are unwell, or if something described in these pages
resembles something that happened to you, that is a matter for a qualified
clinician. Nothing here changes that, and nothing here substitutes for it.
\end{minipage}
\end{center}
\vspace*{\fill}
\cleardoublepage

% --- Epigraph ----------------------------------------------------------
\thispagestyle{empty}
\vspace*{2.6in}
\begin{center}
\begin{minipage}{0.74\textwidth}
\itshape\centering
Every tradition that ever engineered a mind used structure --- repetition,
anchoring, the planted trigger. None of them gave you the schematic. This
book is the schematic, drawn so you can see what has already been built
around you.
\end{minipage}
\end{center}
\clearpage

% --- Table of Contents -------------------------------------------------
{%
  \hypersetup{linkcolor=black}%
  \tableofcontents*
}
\clearpage

% --- List of Figures ---------------------------------------------------
{%
  \hypersetup{linkcolor=black}%
  \listoffigures
}
\clearpage

% --- Preface, How to Use, Quick-Start (unnumbered front matter) --------
\input{front_prose.tex}
